\begin{trapbox}
	\begin{description}[style=nextline, leftmargin=2.8em, labelsep=0.5em, itemsep=0.35em]
		\item[\textbf{\textcolor{CTrap}{易错点一（符号混淆）：}}] 误记 $y_1y_2 = p^2$，忽略了负号．
		\item[\textbf{\textcolor{CTrap}{正确理解（符号来源）：}}] 由韦达定理，$y_1y_2 = -p^2$，负号来自代入消元后的常数项．
		\item[\textbf{\textcolor{CTrap}{错因分析（推导细节）：}}] 代入直线 $x=my+\frac{p}{2}$ 后得 $y^2-2pmy-p^2=0$，常数项为 $-p^2$，故积为负．学生容易遗忘负号，需特别注意．
	\end{description}

	\medskip
	\begin{description}[style=nextline, leftmargin=2.8em, labelsep=0.5em, itemsep=0.35em]
		\item[\textbf{\textcolor{CTrap}{易错点二（形式误用）：}}] 将结论直接套用在 $x^2=2py$ 等非标准形式抛物线上．
		\item[\textbf{\textcolor{CTrap}{正确理解（适用条件）：}}] 结论仅限于 $y^2=2px\ (p>0)$ 成立，其他形式需先转化为标准形式再使用对应结论．
		\item[\textbf{\textcolor{CTrap}{错因分析（形式差异）：}}] 不同标准方程焦点位置不同，代入韦达后常数项符号不同，导致 $y_1y_2$ 公式改变．
	\end{description}

	\medskip
	\begin{description}[style=nextline, leftmargin=2.8em, labelsep=0.5em, itemsep=0.35em]
		\item[\textbf{\textcolor{CTrap}{易错点三（斜率遗漏）：}}] 设点斜式 $y=k(x-\frac{p}{2})$ 联立时，未检验斜率不存在的情况．
		\item[\textbf{\textcolor{CTrap}{正确理解（通径验证）：}}] 当 $k$ 不存在时，弦为通径 $x=\frac{p}{2}$，此时 $x_1 = x_2$，且 $x_1 = \frac{p}{2}$，$y_1=p$，$y_2=-p$，仍满足 $x_1x_2=\frac{p^2}{4},y_1y_2=-p^2$．
		\item[\textbf{\textcolor{CTrap}{错因分析（设线不完整）：}}] 使用 $x=my+\frac{p}{2}$ 可以统一处理所有情形，无需分类讨论，初学者容易遗漏．
	\end{description}
\end{trapbox}
